\textbf{Время:} 01:15

\textbf{Тезис:} общая бозонная симметрия связывает историю оптических и пионных корреляций, но измерительные процедуры различны.

\textbf{Визуал:} три ступени: фотоны и HBT → пионные пары → распределённый источник в ядерном столкновении.

\textbf{Текст:} Пионы, как и фотоны, — бозоны: их спин целый, а квантовое описание симметрично при перестановке одинаковых частиц. В работе Голдхаберов, Ли и Пайса 1960 года рассматривались пионы при аннигиляции антипротона с протоном. Симметризация одинаково заряженных пионов позволяла качественно объяснить наблюдавшиеся угловые корреляции. В модельном расчёте они зависели от размера объёма испускания. Возникала обратная задача: какие размеры совместимы с данными? Это ещё не было точным измерением радиуса тяжёлоионного источника.

Общая основа — свойства тождественных бозонов, то есть частиц с целым спином, и информация в парах. Но телескоп в лабораторию буквально не переносится: в оптике меняли базу приёмников, а здесь сравнивают импульсы двух частиц одного события.

В 1970-е годы Копылов и Подгорецкий развивали связь таких корреляций с пространственными и временными характеристиками испускания. Чтобы воспользоваться этой идеей, сначала научимся отделять совместную регистрацию от простого большого числа частиц.

\textbf{Переход:} Что именно мы будем называть корреляцией и с чем сравнивать число пар?
